
\part{Root constructions and the compact comparison}
\chapter{The Igusa product and its root presentation}
The seed is the weak Jacobi form
\[
 \phiweak(\tau,z)=\sum_{n\ge0,\,\ell\in\Z}f(n,\ell)q^nr^\ell,
 \qquad f(0,0)=10.
\]
The singular theta lift forms the monic weight-five product
\[
 \Dmon(Z)=q^{1/2}r^{1/2}s^{1/2}
 \prod_{(n,\ell,m)>0}(1-q^nr^\ell s^m)^{f(nm,\ell)},
 \qquad \ChiTen=\Dmon^2.
\]
The product records signed multiplicities, a Weyl prefactor, and an
automorphy law. The root bracket is supplied by the
Gritsenko--Nikulin presentation. A geometric orientation and an
operator trace require their own carriers and maps.

\begin{definition}[Igusa root constructions]
Fix the Gritsenko--Nikulin Lie superalgebra \(\gIg\), its positive
subalgebra \(\nIg\), and a proper additive height on its positive
root cone. Its root constructions are the height-complete
Chevalley coalgebra, enveloping algebra, and level-zero current
envelope. Hong's quantum group also uses the full
Borcherds--Cartan superdatum. The diagonal Fredholm model uses a
chosen parity realization, Hilbert norms, and an absolute-convergence
domain. Write \(\RealDatum\) for this collection with those choices.
A map of collections includes a specified map on every component.
\end{definition}

\begin{theorem}[Universal primitive map and PBW recognition]
Let \(\mathfrak p\) be a positive-root-graded Lie superalgebra
over \(\mathbb C\), with finite-dimensional homogeneous pieces.
Suppose a homogeneous parity-preserving map from every simple-root
generating space of \(\nIg\) to \(\mathfrak p\) satisfies
the defining Borcherds--Serre relations. It extends uniquely to
a graded Lie-superalgebra map
\[
 f:\nIg\longrightarrow\mathfrak p.
\]
It induces continuous maps on height completions, enveloping
algebras, Chevalley chains, and level-zero currents. The map \(f\)
is an isomorphism if and only if it is surjective and the root-graded
PBW series agree in each parity.
\end{theorem}
\begin{proof}
Extend the chosen maps to the free Lie superalgebra on the
simple-root spaces. The relation hypothesis makes the defining
ideal lie in its kernel, giving \(f\). The generators prove
uniqueness. On enveloping algebras the map sends
\(x_1\cdots x_j\) to \(f(x_1)\cdots f(x_j)\).
On Chevalley chains it is the graded symmetric map on the shifted
generators. It commutes with the Chevalley differential because
\(f\) preserves brackets. The current map is
\(\partial^j x\mapsto\partial^j f(x)\), preserving
\([x_\lambda y]=[x,y]\). Root degree preservation gives
compatible maps at each height and hence on the stated limits.

If \(f\) is surjective, \(U(f)\) is surjective. Proper
positive height makes each root-graded PBW component
finite-dimensional. Equality of its even and odd dimensions makes
\(U(f)\) bijective in every component. PBW embeds each Lie
superalgebra into its enveloping algebra, so \(f\) is injective.
The converse follows by applying the constructions to its inverse.
\end{proof}


A Hall-category comparison must intertwine the actual extension
correspondences. A quantum-double comparison preserves its coproduct
and pairing. A Fredholm comparison intertwines the operators on
their declared domains. These maps are part of a realization of
\(\RealDatum\). The universal primitive map supplies the algebraic
component from which they are to be constructed.

\begin{theorem}[Integral Igusa identity]
The stated root constructions give the identities
\[
 \den(\gIg)=\Dmon(2Z),
 \qquad
 e^{-\rho}\sch\widehat C_*(\nIg)=\Dmon(2Z),
\]
\[
 \operatorname{sdet}_F(1-K_Z)=e^{\rho}\Dmon(2Z),
 \qquad
 \mathcal Z^{\red}_{K3\times E}=-\ChiTen^{-1}.
\]
After the half-coordinate pullback $\jmath(Z)=Z/2$, the reduced scalar is the inverse square of the first-order denominator:
\[
 \mathcal Z^{\red}_{K3\times E}
 =-\bigl(\jmath^*\den(\gIg)\bigr)^{-2}.
\]
\end{theorem}
\begin{proof}
The first identity is the Gritsenko--Nikulin denominator theorem.  PBW identifies the second with the Euler supercharacter of Chevalley chains.  Trace-class convergence gives the third as the product of root eigenvalues.  The fourth is the primitive reduced curve-counting theorem of Oberdieck--Pixton in monic normalization.  Substitution of the half-coordinate identity gives the final equality.
\end{proof}

\chapter{The compact geometric refinement}
The physical charge of a brane is additive.  The automorphic Gram degree is quadratic.  Their exact relation is carried by the polarization cocycle
\[
 B(c_1,c_2)=\PiGram(c_1+c_2)-\PiGram(c_1)-\PiGram(c_2).
\]
The central extension
\[
 (c,\lambda)(d,\mu)=(c+d,\lambda+\mu+B(c,d))
\]
makes $c\mapsto(c,\PiGram(c))$ additive.  This extension is the natural grading object for a Hall source whose character descends to the three Igusa variables.

\begin{definition}[Geometric realization datum]
A geometric $K3\times E$ realization consists of a reduced oriented extension theory on compactly supported objects, its additive charge grading, the polarized Gram extension, projection-finite and wrapped elliptic sectors, mixed extension correspondences, a primitive integration map to the universal root datum, a coproduct and Hopf pairing, and an orientation-line map to $\mathcal L^5\tensor\nu_5$.
\end{definition}

\begin{recognition}[Geometric Hall--Borcherds recognition]
Let the geometric realization datum supply a positive-root-graded
primitive Lie superalgebra \(\mathfrak p_X\) and simple geometric
classes satisfying the defining relations of \(\nIg\).
The resulting map \(\nIg\to\mathfrak p_X\) is an isomorphism
precisely when it is surjective and the two parity-refined PBW
series agree. It then identifies their completed enveloping
algebras, Chevalley chains, and level-zero current envelopes.
\end{recognition}
\begin{proof}
Apply the universal primitive map and PBW recognition theorem to
these geometric classes. The induced maps are the tensor-word,
Chevalley, and current maps constructed in its proof.
\end{proof}

The compact construction must produce these primitive classes from
oriented reduced coefficients on extension stacks. It must prove
their relations and generation and identify both parity dimensions.
To compare Hall doubles, centres, or Pfaffian sections, it must also
construct the coproduct, monoidal, and line maps. These maps have a
common geometric source. Their existence does not follow from the
PBW character.
